\documentclass[a4paper,14pt]{extarticle}

\usepackage[T2A]{fontenc}
\usepackage[utf8]{inputenc}
\usepackage[russian]{babel}
\usepackage{geometry}
\usepackage{setspace}
\usepackage{indentfirst}
\usepackage{array}
\usepackage{booktabs}
\usepackage{listings}
\usepackage{xcolor}

\geometry{
    left=30mm,
    right=15mm,
    top=20mm,
    bottom=20mm
}

\onehalfspacing
\setlength{\parindent}{1.25cm}
\setlength{\parskip}{0pt}
\setlength{\emergencystretch}{3em}

\definecolor{codebg}{RGB}{248,248,248}
\lstdefinestyle{cudastyle}{
    language=C++,
    basicstyle=\ttfamily\normalsize,
    frame=single,
    breaklines=true,
    showstringspaces=false
}

\begin{document}

\begin{titlepage}
\begin{center}
\textbf{МОСКОВСКИЙ АВИАЦИОННЫЙ ИНСТИТУТ}\\
\textbf{(НАЦИОНАЛЬНЫЙ ИССЛЕДОВАТЕЛЬСКИЙ УНИВЕРСИТЕТ)}\\
Институт №8 «Компьютерные науки и прикладная математика»
\end{center}

\vspace{3.5cm}
\begin{center}
\textbf{\Large Лабораторная работа №5}\\
\vspace{0.3cm}
по курсу «Технологии параллельного программирования»\\
\vspace{0.8cm}
\textbf{Сортировка чисел на GPU}
\end{center}

\vspace{4cm}
\begin{flushright}
Выполнил: Ахметшин Б.Р.\\
Группа: М8О-403Б-22\\
Преподаватели: А.Ю. Морозов,\\
\hspace*{3.6cm} Е.Е. Заяц
\end{flushright}

\vfill
\begin{center}
Москва, 2026
\end{center}
\end{titlepage}

\section*{Условие}
Цель работы: реализовать сортировку чисел на GPU с использованием
фундаментальных примитивов (гистограмма, сканирование), провести сравнение
с CPU и исследовать производительность с помощью \texttt{nvprof}.

Выполненная реализация соответствует \textbf{варианту 3}:
сортировка подсчетом для чисел типа \texttt{uchar} (диапазон 0..255).

Формат данных:
\begin{itemize}
    \item вход: \texttt{[int32 n][n x uchar]}, чтение из \texttt{stdin};
    \item выход: \texttt{n} отсортированных \texttt{uchar} в \texttt{stdout}.
\end{itemize}

\section*{Программное и аппаратное обеспечение}
\textbf{GPU}

\begin{tabular}{>{\raggedright\arraybackslash}p{0.50\textwidth}>{\raggedright\arraybackslash}p{0.34\textwidth}}
Имя & NVIDIA Tesla T4 \\
Compute Capability & 7.5 \\
Объем глобальной памяти & 15637086208 \\
Разделяемая память на блок & 49152 \\
Константная память & 65536 \\
Регистров на блок & 65536 \\
Макс. потоков на блок & (1024, 1024, 64) \\
Макс. размер сетки & (2147483647, 65535, 65535) \\
Количество мультипроцессоров & 40 \\
\end{tabular}

\vspace{0.4cm}
\textbf{CPU}

\begin{tabular}{>{\raggedright\arraybackslash}p{0.50\textwidth}>{\raggedright\arraybackslash}p{0.34\textwidth}}
Архитектура & x86\_64 \\
CPU op-mode(s) & 32-bit, 64-bit \\
Address sizes & 48 bits physical, 48 bits virtual \\
Порядок байт & Little Endian \\
CPU(s) & 12 \\
Имя модели & AMD Ryzen 5 7600 6-Core Processor \\
Ядер / потоков & 6 / 12 \\
\end{tabular}

\vspace{0.4cm}
\textbf{RAM}\\
64 GB

\vspace{0.2cm}
\textbf{ОС}\\
Ubuntu 24.04.4 LTS

\vspace{0.2cm}
\textbf{Компилятор}\\
nvcc, g++

\section*{Метод решения}
Используется сортировка подсчетом для \texttt{uchar}:
\begin{enumerate}
    \item строится гистограмма частот значений 0..255 с атомарными операциями;
    \item выполняется exclusive scan (Blelloch) по 256 корзинам
    с бесконфликтным доступом к shared memory;
    \item по \texttt{histogram} и \texttt{offsets} восстанавливается
    отсортированный массив.
\end{enumerate}

На CPU реализован эквивалентный counting sort для корректного сравнения
времени и верификации результата.

\section*{Описание программы}
\begin{itemize}
    \item \texttt{lab5.cu} --- GPU-реализация (гистограмма + scan + восстановление массива),
    внутренние замеры \texttt{kernel\_ms}/\texttt{total\_ms}.
    \item \texttt{lab5-cpu.cpp} --- CPU-реализация counting sort,
    внутренний замер \texttt{total\_ms}.
    \item \texttt{generate\_case.py} --- генератор бинарных тестов
    формата \texttt{[int32 n][n x uchar]}.
    \item \texttt{run\_case.sh} --- запуск исполняемого файла на входном кейсе.
    \item \texttt{profile\_nvprof.sh} --- запуск с профилировщиком
    \texttt{nvprof} (сводный отчет).
\end{itemize}

Параметры через переменные окружения:
\begin{itemize}
    \item GPU: \texttt{LAB5\_HIST\_THREADS};
    \item GPU: \texttt{LAB5\_FILL\_THREADS};
\end{itemize}

\textbf{Launch-конфигурации kernel'ов из кастомных параметров:}
\begin{itemize}
    \item \texttt{clearHistogramKernel}:\\
    \texttt{block = (LAB5\_HIST\_THREADS, 1, 1)},\\
    \texttt{grid.x = ceil(256/LAB5\_HIST\_THREADS)}.
    \item \texttt{buildHistogramKernel}:\\
    \texttt{block = (LAB5\_HIST\_THREADS, 1, 1)},\\
    \texttt{grid.x = ceil(n/LAB5\_HIST\_THREADS)}.
    \item \texttt{scan256Kernel}:\\
    \texttt{block = (128, 1, 1)}, \texttt{grid = (1, 1, 1)} (фиксировано для 256 корзин).
    \item \texttt{fillSortedArrayKernel}:\\
    \texttt{block = (LAB5\_FILL\_THREADS, 1, 1)},\\
    \texttt{grid = (256, 1, 1)}.
\end{itemize}

\section*{Результаты}
\textbf{Замеры времени:}

\begin{lstlisting}[basicstyle=\ttfamily\scriptsize,frame=single]
+----+-----------+---------------------+-----------+----------------+---------------+
| No | n         | CUDA config         | CPU, ms   | GPU kernel, ms | GPU total, ms |
+----+-----------+---------------------+-----------+----------------+---------------+
| 1  | 1000      | hist=64, fill=32    | 0.002250  | 0.323168       | 395.620671    |
+----+-----------+---------------------+-----------+----------------+---------------+
| 2  | 1000      | hist=128, fill=64   | 0.002250  | 0.317088       | 293.304433    |
+----+-----------+---------------------+-----------+----------------+---------------+
| 3  | 1000      | hist=256, fill=128  | 0.002250  | 0.352256       | 290.412191    |
+----+-----------+---------------------+-----------+----------------+---------------+
| 4  | 1000000   | hist=64, fill=32    | 0.955771  | 0.274272       | 282.703727    |
+----+-----------+---------------------+-----------+----------------+---------------+
| 5  | 1000000   | hist=128, fill=62   | 0.955771  | 0.284224       | 282.320927    |
+----+-----------+---------------------+-----------+----------------+---------------+
| 6  | 1000000   | hist=256, fill=128  | 0.955771  | 0.296992       | 271.056927    |
+----+-----------+---------------------+-----------+----------------+---------------+
| 7  | 100000000 | hist=64, fill=32    | 92.699408 | 2.894880       | 532.688073    |
+----+-----------+---------------------+-----------+----------------+---------------+
| 8  | 100000000 | hist=128, fill=62   | 92.699408 | 2.714176       | 460.864727    |
+----+-----------+---------------------+-----------+----------------+---------------+
| 9  | 100000000 | hist=256, fill=128  | 92.699408 | 2.991680       | 470.329608    |
+----+-----------+---------------------+-----------+----------------+---------------+
\end{lstlisting}

\textbf{Логи профилировщика \texttt{nvprof}:}

\begin{lstlisting}[basicstyle=\ttfamily\scriptsize,frame=single,breaklines=true]
==10076== NVPROF is profiling process 10076, command: ./lab5
==10076== Profiling application: ./lab5
==10076== Profiling result:
            Type  Time(%)      Time     Calls       Avg       Min       Max  Name
 GPU activities:   48.13%  21.810ms         1  21.810ms  21.810ms  21.810ms  [CUDA memcpy DtoH]
                   46.06%  20.870ms         1  20.870ms  20.870ms  20.870ms  [CUDA memcpy HtoD]
                    3.61%  1.6347ms         1  1.6347ms  1.6347ms  1.6347ms  buildHistogramKernel(unsigned char const *, int, unsigned int*)
                    2.19%  991.12us         1  991.12us  991.12us  991.12us  fillSortedArrayKernel(unsigned int const *, unsigned int const *, unsigned char*)
                    0.01%  6.2080us         1  6.2080us  6.2080us  6.2080us  scan256Kernel(unsigned int const *, unsigned int*)
                    0.01%  3.2000us         1  3.2000us  3.2000us  3.2000us  clearHistogramKernel(unsigned int*)
\end{lstlisting}

\textbf{Анализ профилировщика:}
\begin{itemize}
    \item По сводке \texttt{nvprof} основной вклад в GPU-время дают копирования:
    \texttt{[CUDA memcpy DtoH]} --- 48.13\%, \texttt{[CUDA memcpy HtoD]} --- 46.06\%.
    Суммарно это около 94.19\% времени GPU-активностей.
    \item Сами вычислительные kernel'ы занимают малую долю:
    buildHistogram --- 3.61\%,
    fillSortedArray --- 2.19\%,
    scan256/clearHistogram --- менее 0.02\%.
    \item По таблице замеров видно, что время GPU-ядра мало
    (от 0.27 до 2.99 ms) и растет с \(n\) значительно медленнее, чем время CPU.
    Однако полное время GPU (\(271\!-\!533\) ms) остается высоким из-за накладных
    расходов на обмен данными и инфраструктурные операции запуска.
    \item Влияние конфигурации \texttt{hist/fill} заметно, но умеренно:
    наиболее удачные комбинации уменьшают общее время на десятки миллисекунд,
    не меняя качественно вывод о доминировании накладных расходов.
    \item Возможные улучшения: закрепленная (pinned) host-память, повторное
    использование буферов и контекста между запусками, обработка данных
    пакетами без возврата на CPU после каждого шага.
\end{itemize}

\section*{Выводы}
По результатам экспериментов получены следующие выводы:
\begin{itemize}
    \item Для всех исследованных размеров (\(n=10^3, 10^6, 10^8\))
    текущая GPU-реализация проигрывает CPU по полному времени:
    минимум 290.41 ms против 0.00225 ms при \(n=10^3\),
    271.06 ms против 0.95577 ms при \(n=10^6\),
    460.86 ms против 92.70 ms при \(n=10^8\).
    \item При этом по времени самих kernel'ов GPU работает эффективно:
    0.27--2.99 ms, что существенно меньше CPU-вычислений на больших \(n\).
    \item Профилирование подтверждает, что основной узкий участок --- не вычисления,
    а передача данных HtoD/DtoH, которая занимает более 94\% GPU-активностей.
    \item Лучшая конфигурация из проверенных:
    для \(n=10^3\) и \(10^6\) --- \texttt{hist=256, fill=128},
    для \(n=10^8\) --- \texttt{hist=128, fill=62}.
    \item Итог: алгоритм на GPU корректен и быстрый на уровне kernel'ов, но для
    единичного запуска с копированием данных туда-обратно ускорение не достигается;
    для практического выигрыша нужно сокращать накладные расходы обмена и запуска.
\end{itemize}

\end{document}
